%! TeX program = lualatex
\documentclass[a4paper, 12pt]{article}
% \documentclass[twoside, a4paper, 12pt]{book}

\title{Wstęp do analizy harmonicznej na grupach}
\author{\color{subtext1}Weronika Jakimowicz}
\date{Lato 2025/26}

\usepackage[pl]{../../template}

\usepackage{halloweenmath}

\let\oldamalg\amalg
\newcommand{\uu}[1]{\underline{#1}}
\newcommand{\Mm}{\ensuremath{\mathfrak{M}}}
\newcommand{\lr}{\leftrightarrow}

\DeclareMathOperator{\Cn}{Cn}
\DeclareMathOperator{\Th}{Th}

\DeclareMathOperator{\Hh}{\mathcal{H}}
\DeclareMathOperator{\U}{\mathcal{U}}
\DeclareMathOperator{\B}{\mathcal{B}}


% \usepackage{yhmath} 
% \renewcommand{\Hat}[1]{\widehat{#1}}
%
% \newcommand\widecheck[1]{%
%     #1%
%     \text{\llap{%
%         \smash{\raisebox{\dimexpr 2\depth-0.5ex\relax}{%
%             \scalebox{1}[-1]{$\widehat{\phantom{#1}}$}%
%         }}%
%     }\rule{0pt}{1.1em}}%
% }
%
% \newcommand{\taH}[1]{\widecheck{#1}}
\DeclareFontFamily{U}{mathx}{}
\DeclareFontShape{U}{mathx}{m}{n}{<-> mathx10}{}
\DeclareSymbolFont{mathx}{U}{mathx}{m}{n}
\DeclareMathAccent{\widehat}{0}{mathx}{"70}
\DeclareMathAccent{\widecheck}{0}{mathx}{"71}
\renewcommand{\Hat}[1]{\widehat{#1}}
\newcommand{\taH}[1]{\widecheck{#1}}

\newtcolorbox{zadanie}[2]{
  theoremBoxTHM={green}{green},
    title={Lista #1: zadanie #2}
}

\begin{document}
% \frontmatter 
\maketitle
\thispagestyle{empty}
\setcounter{page}{0}

\tableofcontents

\pagestyle{fancy}

\section{C${}^*$ algebra basics}

\begin{definition}{}{}
  \acc{Algebra Banacha} $A$ na której jest zdefiniowana operacja ${}^*$ taka, że
  \begin{enumerate}
    \item $(x+y)^*=x^*+y^*$
    \item $(\lambda x)^*=\overline{\lambda}x^*$
    \item $(xy)^*=y^*x^*$
    \item $x^{**}=x$
    \item $\|x^*x\| = \|x\|^2$
  \end{enumerate}
  jest nazywana \buff{$C^*$-algebrą}
\end{definition}

$l^1\Z$ nie jest $C^*$-algebrą, bo chyba mogę wziąć $x_n=e^{ni\frac{2\pi}{k}}$, wtedy 
$$x_n^*x_n=(e^{-ni\frac{2\pi}{k}})(e^{ni\frac{2\pi}{k}})$$
ma na $n$-tej współrzędnej
$$\sum_{m=0}^n e^{-mi\frac{2\pi}{k}}e^{(n-m)i\frac{2\pi}{k}}=e^{ni\frac{2\pi}{k}}\sum e^{-2mi\frac{2\pi}{k}}$$










\subsection{Gelfand theory}

\begin{definition}{spektrum algebry}{}
  Zbiór wszystkich multiplikatywnych niezerowych funkcjonałów na $A$ (czyli homomorfizmów $A\to \C$) jest nazywany \buff{spektrum algebry}, $\sigma(A)$.
\end{definition}

Na $\sigma(A)$ dajemy topologię *-słabą, czyli zbieżności punktowej na $A$.

$\sigma(A)$ jest domkniętym podzbiorem kuli jednostkowej w $A^*$, bo
$$\sigma(A)=\{h\in A^*\;:\;|h|\leq 1\;\land\;h(1)=1\;\land\;(\forall\;x,y\in A)\;h(xy)=h(x)h(y)\}$$
a warunek $h(xy)=h(x)h(y)$ jest zachowywany przez granice punktowe.

\begin{theorem}{}{}
  Domknięcie $\overline{I}$ właściwego ideału $I\subseteq A$ jest ideałem właściwym.
\end{theorem}

\begin{proof}
  Zbiór elementów odwracalnych jest otwarty.
\end{proof}

\begin{theorem}{}{}
  Niech $A$ będzie przemienną algebrą Banacha. Wtedy $h\mapsto \ker h$ to bijekcja pomiędzy $\sigma(A)$ i maksymalnym spektrum $A$
\end{theorem}

\begin{proof}
  $h\mapsto \ker h$ injekcja, bo
  \begin{itemize}
    \item dla $h\in\sigma(A)$ mamy $h(1)=1\neq 0$ więc $\ker h$ ma kowymiar $1$
    \item $\ker h=\ker g$ oznacza, że $h=g$ bo $x=h(x)+y$ dla $y\in \ker h$, więc $g(x)=h(x)g(1)+g(y)=h(x)$
  \end{itemize}

  W drugą stronę, niech $\mathfrak{m}$ będzie ideałem maksymalnym, wtedy
  \begin{itemize}
    \item $A/\mathfrak{m}$ jest algebrą Banacha bez nietrywialnych ideałów - wszystko odwracalne
    \item Gelfand-Mazur: $A/\mathfrak{m}\cong\C$
    \item $\ker(A\to A/\mathfrak{m}\to \C)=\mathfrak{m}$
  \end{itemize}
\end{proof}

\begin{definition}{transformata Gelfanda}{}
  Niech $x\in A$, wtedy $\Hat{x}$ to ciągła funkcja na $\sigma(A)$:
  $$\Hat{x}(h)=h(x).$$

  Odwzorowanie $x\mapsto \Hat{x}$ to \buff{transformata Gelfanda} na $A$.
\end{definition}

\begin{theorem}{}{}
  Niech $A$ będzie przemienną algebrą Banacha z jedynką. Wtedy 
  \begin{enumerate}
    \item $x$ jest odwracalne $\iff$ $\Hat{x}$ nigdy nie znika
    \item $\im(\Hat{x})=\sigma(x)$
    \item $\|x\|_\infty=\rho(x)\leq\|x\|$
  \end{enumerate}
\end{theorem}

\begin{proof}
  \begin{enumerate}
    \item $x$ nie jest odwracalne $\iff$ ideał generowany przez $x$ jest właściwy $\iff$ $x$ jest zawarty w pewnym ideale maksymalnym w $A$ $\iff$ $x\in \ker h$ dla pewnego $h$ $\iff$ $\Hat{x}$ ma zero
    \item $\lambda \in\sigma(x)$ $\iff$ $x-\lambda$ jest nieodwracalne $\iff$ $\Hat{x}(h)-\lambda=0$ dla jakiś $h\in \sigma(A)$
    \item \color{red}wynika wprost z poprzedniego
  \end{enumerate}
\end{proof}











\section{Grupa dualna}

$$L_yf(x)=f(y^{-1}x)$$

\begin{theorem}{}{}
  Nieredukowalne reprezentacje lokalnie zwartej abelowej grupy są jednowymiarowe.
\end{theorem}

Czyli każda reprezentacja lokalnie zwartej grupy abelowej $G$ działa na $\C$ i możemy ją zapisać jako $\pi(x)z=\xi(x)z$ dla $\xi:G\to S^1$ (\buff{charakter}).

\begin{definition}{positive type function}{}
Funkcja dodatnio określona (positive type function) to funkcja $\phi\in L^\infty G$ taka, że dla wszystkich $f\in L^1G$ 
$$\int (f^*\ast f)\phi\geq 0$$
\end{definition}

$\Hat{G}$ to zbiór wszystkich charakterów $G$, jest on zawarty w zbiorze funkcji dodatnio określonych o normie jeden. 

\begin{theorem}{}{}
  Dla każdego $\xi\in\Hat{G}$ definiujemy niezdegenerowaną *-reprezentację $L^1G$
  $$\xi(f)=\int \langle x,\xi\rangle f(x)dx=\int\xi(x)f(x)dx$$
  To daje utożsamienie $\Hat(G)$ ze spektrum $L^1G$.
\end{theorem}

\subsection{Topologia na $\Hat{G}$}

\begin{definition}{}{}
  Topologia na $\Hat{G}$ to topologia zbieżności na zbiorach zwartych w $G$. 
\end{definition}

Inna nazwa: topologia jednostajnej zbieżności na zbiorach zwartych w $G$. Baza topologii to zbiory 
$$N(\phi_0, \epsilon, K)=\{\phi:(\forall\;x\in K)\;|\phi(x)-\phi_0(x)|<\epsilon\}$$
dla $\phi_0$ w $\Hat{G}$, $\epsilon>0$ i $K\subseteq G$ zwartego.


\begin{definition}{*-słaba topologia}{}
  To topologia na $X$ w której wszystkie funkcje $\langle x, -\rangle:X^*\to \C$, $\langle x,\phi\rangle =\phi(x)$, są ciągłe.
\end{definition}

Czyli na $L^\infty G$ zbiory otwarte to $\{f\in L^\infty\;:\;\int fg\in U\}$ dla $g\in L^1G$ i $U\subseteq \C$ otwartego.

\begin{theorem}{}{}
  Topologia zwartej zbieżności na $\Hat{G}$ pokrywa się z *-słabą topologią dziedziczoną z $L^\infty G$.
\end{theorem}

\begin{proof}\color{red}
  3.31 z follanda, na razie nie mam czasu
\end{proof}

\begin{zadanie}{2}{2}
  Niech $G$ będzie grupą zwartą, niech $\chi:G\to U(1)$ będzie homomorfizmem. Zauważ, że $\chi\in L^p(G)$ dla każdego $p$.
\end{zadanie}

\begin{zadanie}{2}{3}
  Przy założeniach poprzedniego zadania pokaż, że jeśli $\chi$ nie jest stały, to $\int_G\chi=0$. Wskazówka: użyj niezmienniczości miary Haara.
\end{zadanie}

\begin{proof}
  Nie wprost: weźmy $h\in G$ takie, że $\chi(h^{-1})\neq 1$. Wtenczas 
  $$\int_G\chi(g)dg=\int_GL_h\chi(g)dg=\int_G\chi(h^{-1}g)dg=\int_G\chi(h^{-1})\chi(g)dg=\chi(h^{-1})\int_G\chi(g)dg$$
  czyli
  $$0=(1-\chi(h^{-1}))\int_G\chi(g)dg$$
  no a skoro $\chi(h^{-1})\neq 1$ to musowo $\int\chi=0$.
\end{proof}

\begin{zadanie}{2}{4}\label{2:4}
  Wywnioskuj, że dla $\chi,\tau\in\Hat{G}$, $\chi=\tau$ albo $\chi$ i $\tau$ są ortogonalne w $L^2G$
\end{zadanie}

\begin{proof}
  $$\langle \chi, \tau\rangle = \int\chi\overline{\tau}$$
  $\chi\overline{\tau}\in\Hat{G}$, więc mamy dwie opcje: 
  \begin{itemize}
    \item albo $\int \chi\overline{\tau}\neq 0$, więc $\chi\overline{\tau}$ jest stale równe 1, to znaczy że $\overline{\tau}=\overline{\chi}$ czyli $\chi=\tau$, 
    \item albo $\int \chi\overline{\tau}=0$ czyli charaktery są ortogonalne.
  \end{itemize}
\end{proof}

\begin{zadanie}{2}{5}
  Jeśli $G$ jest zwarta, to $\Hat{G}$ jest dyskretna. Wskazówka: zauważ, że $U=\{f\in L^\infty G\;:\;\int f>\frac{1}{2}\}$ jest otwarty (w *-słabej topologii). Pokaż, że jeśli $\chi\in \Hat{G}\cap U$ to $\chi=1$.
\end{zadanie}

\begin{proof}
  Skoro $G$ jest zwarta, to funkcja stale równa $1$ jest w $L^1G$. W takim razie 
  $$\{f\in L^\infty G\;:\;\left|\int fd\lambda\right|=\left|\int f\cdot 1d\lambda\right|>\frac{1}{2}\}$$
  jest zbiorem otwartym w *-słabej topologii. Z zadania \ref{2:4} dla dowolnego $\chi \in \Hat{G}$ jest $\int \chi=1$ gdy $\chi=1$ i $\int\chi=0$ w przeciwny przypadku. Stąd $\{1\}$ jest zbiorem otwartym w $\Hat{G}$, więc każdy singleton też jest zbiorem otwartym.
\end{proof}

\begin{theorem}{}{}
  Jeśli $G$ jest dyskretna, to $\Hat{G}$ jest zwarta.
\end{theorem}

\begin{proof}
  $G$ dyskretna $\implies$ $L^1G$ ma jednostkę, czyli funkcję $\delta$ która jest $1$ w jedynce i $0$ na całej reszcie. Jeśli $\sigma(G)\cong \Hat{G}$ jest niezwarta, to $\Hat{x}$, definiowane jako $\Hat{x}(h)=h(x)$ dla $h\in\sigma(G)$, znika w nieskończoności dla każdego $x\in G$ (przed 1.30 z Follanda). {\color{red}dokończyć}
\end{proof}

\subsection{Liczby p-adyczne}

\begin{zadanie}{2}{8}
  Niech $\Q_p$ oznacza grupę addytywną ciała $p$-adycznego. Każdy $x\in\Q_p$ jest reprezentowany przez sumę $\sum_{-\infty \leq k\leq \infty} c_kp^k$ gdzie $c_k\in (0,...,p-1)$ oraz istnieje takie $N_0\in\Z$, że $c_k=0$ jeśli $k<N_0$. Zdefiniujmy 
  $$\chi_1(x)=\exp\left(2\pi i\sum_{k\leq -1} c_kp^k\right).$$
  Pokaż, że $\chi_1\in\Hat{\Q_p}$. Ogólniej, dla $y\in\Q_p$ zdefiniujmy $\chi_y(x)=\chi_1(xy)$. Pokaż, że $\chi_y\in\Hat{\Q_p}$.
\end{zadanie}

\begin{zadanie}{2}{9}
  Pokaż, że $\chi_y\chi_{y'}=\chi_{y+y'}$.
\end{zadanie}

\begin{proof}
\end{proof}















\subsection{Transformata Fouriera}

$\Hat{G}$ identyfikujemy z $\sigma(L^1G)$: dla $\chi\in\Hat{G}$ przypisujemy
$$f\mapsto \overline{\chi}(f)=\chi^{-1}(f)=\int\overline{\langle x,\chi\rangle}f(x)dx$$

Transformata Gelfanda=Fouriera na $L^1G$ to wtenczas 
$$L^1G\ni f\mapsto \Hat{f}(\xi)=\int\overline{\langle x,\xi\rangle}f(x)dx\in C(\Hat{G})$$

\begin{zadanie}{3}{1}
  Niech $f\in C_c^\infty(\R)$. Oblicz $\Hat{-if}'$.
\end{zadanie}

\begin{proof}
  Czym jest $\Hat{-if}$? 
  \begin{align*}
    \Hat{-if}(\xi)&=\int \overline{\langle x,\xi\rangle}(-if(x))dx = \\ 
                  &=\int \overline{\langle x, \xi\rangle i}f(x)dx
  \end{align*}
  ale $\Hat{\R}\cong \R$ przez $\langle x,\xi\rangle=e^{2\pi i\xi x\rangle}$. Korzystając z tego liczymy, że pochodna względem $\xi$ to (według reguły Leibniza) powinno być
  \begin{align*}
    \frac{d}{d\xi}\Hat{-if}(\xi)&=\int \frac{\partial}{\partial \xi}\overline{i\langle x,\xi\rangle}f(x)dx=\\ 
                                &=\int \frac{\partial}{\partial\xi}\overline{i e^{2\pi i \xi x}}f(x)dx= \\ 
                                &=\int \overline{-2\pi x e^{2\pi i \xi x}}f(x)dx=-2\pi\Hat{xf}(\xi).
  \end{align*}
\end{proof}

\begin{zadanie}{3}{2}
  Niech $f\in L^1A$, $\omega\in\Hat{A}$. Pokaż, że $\Hat{\omega\cdot f}=L_\omega \Hat{f}$.
\end{zadanie}

\begin{proof}
  Rozpisujemy od lewej strony (dokładniej niż niżej)
  \begin{align*}
    \Hat{\omega\cdot f}(\xi)&=\int \overline{\langle x,\xi\rangle}\langle x,\omega\rangle f(x)dx=\\ 
                       &=\int \overline{\langle x,\xi\rangle\langle x,\omega^{-1}\rangle}f(x)dx=\\ 
                       &=\int \overline{\langle x,\xi\omega^{-1}\rangle}f(x)dx=\\ 
                       &=\Hat{f}(\xi\omega^{-1})=\Hat{f}(\omega^{-1}\xi)=L_{\omega}\Hat{f}(\xi)
  \end{align*}
\end{proof}

\begin{theorem}{}{}
  Transformata Fouriera to zwężający *-homomorfizm $L^1G\to C_0(\Hat{G})$ (lub $C(\Hat{G})$ gdy $\Hat{G}$ jest zwarte) której obraz jest gęstym podzbiorem $C_0(\Hat{G})$.
\end{theorem}

To jest bardziej abstrakcyjne postawienie lematu Riemanna-Lesbegue'a

Kilka ważnych własności transformaty Fouriera:
$$\Hat{L_yg}(\xi)=\int\overline{\langle x,\xi\rangle}f(y^{-1}x)dx=\int\overline{\langle yx,\xi\rangle}f(x)dx=\overline{\langle y,\xi\rangle}\Hat{f}(\xi)$$
$$\Hat{(\eta f)}(\xi)=\int\overline{\langle x,\xi\rangle}\langle x,\eta\rangle f(x)dx=\Hat{f}(\eta^{-1}\xi)=L_\eta \Hat{f}(\xi)$$

Dla miary $\mu\in M(G)$ definiujemy 
$$\Hat{\mu}(\xi)=\int\overline{\langle x,\xi\rangle}d\mu(x).$$

\begin{theorem}{}{}
  Odwzorowanie $M(\Hat{G})\ni\mu\to \phi_\mu\in C_b(G)$ definiowane jako
  $$\phi_\mu(x)=\int\langle x,\xi\rangle d\mu(\xi)$$
  jest zwężającą iniekcją.
\end{theorem}

\begin{theorem}{Bochnera}{}
  Jeśli $\phi\in\mathcal{P}(G)$ to istnieje jedyna dodatnia miara $\mu\in M(\Hat{G})$ taka, że $\phi=\phi_\mu$.

  \color{red}Gal robi dla niedodatnich tylko tak ogólniej?
\end{theorem}

\begin{proof}
  Niżej
\end{proof}

Definiujemy 
$$B(G)=\{\phi_\mu\;:\;\mu\in M(\Hat{G})\}$$
oraz $B^p(G)=B(G)\cap L^pG$. Z twierdzenia Bochnera $span\left(\mathcal{P}(G)\right)=B(G)$.

\begin{lemma}{}{}
  Jeśli $K\subseteq\Hat{G}$ jest zwarty to istnieje $f\in C_c(G)\cap \mathcal{P}$ takie, że $\Hat{f}\geq 0$ na $\Hat{G}$ oraz $\Hat{f}>0$ na $K$.
\end{lemma}

Odwzorowanie $\mu\mapsto \phi_\mu$ to bijekcja $M(\Hat{G})\to B(G)$. Jej odwrotność oznaczymy $\phi\mapsto \mu_\phi$.

\begin{theorem}{}{}
  Dla $f,g\in B^1(G)$ mamy $\Hat{f}d\mu_g=\Hat{g}d\mu_f$.
\end{theorem}

\begin{theorem}{Fourier inversion theorem I}{}
  \begin{itemize}
    \item $f\in B^1(G)\implies \Hat{f}\in L^1(\Hat{G})$
    \item miara Haara $d\xi$ na $\Hat{G}$ jest znormalizowana względem miary Haara $dx$ na $G$ to $d\mu_f(\xi)=\Hat{f}(\xi)d\xi$, czyli
      $$f(x)=\int\langle x,\xi\rangle \Hat{f}(\xi)d\xi$$
  \end{itemize}
\end{theorem}

Jeśli $dx$ to miara Haara na $G$, to $d\xi$ z założeń twierdzenia wyżej jest miara Haara na $\Hat{G}$. Nazywamy ją miarą dualna do $dx$.

\begin{zadanie}{3}{7}
  Jeśli $G$ jest zwarta z probabilistyczną miarą Haara, to dualna miara Haara jest miarą liczącą.
\end{zadanie}

\begin{proof}
  Niech $g(x)=1$ będzie funkcją stale równą $1$. Wtedy $\Hat{g}=\chi_{\{1\}}$, bo 
  $$\Hat{g}(\xi)=\int \overline{\langle x, \xi\rangle}g(x)dx=\int \overline{\langle x,\xi\rangle}dx$$
  a skoro charaktery tworzą bazę ortonormalną, to całka po prawej jest $0$ gdy $\xi\neq 1$ i wynosi $1$ dla charakteru trywialnego, czyli stale równego $1$. 

  Stąd 
  $$1=g(x)=\int \langle x, \xi\rangle \Hat{g}(\xi)d\xi=\int\langle x,\xi\rangle \chi_{\{1\}}(\xi)d\xi=d\xi(\{1\})$$
  czyli miara singletona $1$ wynosi $1$, całą resztę osiągamy przesuwaniem bo już pokazaliśmy, że dla zwartego $G$ grupa dualna jest dyskretna.
\end{proof}

\begin{zadanie}{3}{8}
  Jeśli $G$ jest dyskretna z liczącą miarą Haara to dualna miara Haara na $\Hat{G}$ jest miarą probabilistyczną.
\end{zadanie}

\begin{proof}
  Niech $g(x)=\chi_{\{1\}}(x)$ będzie funkcją, wtedy $\Hat{g}(\xi)=1$ z analogicznego powodu co wcześniej:
  $$\Hat{g}(\xi)=\int \overline{\langle x,\xi\rangle}g(x)dx=\int \overline{\langle x, \xi\rangle}\chi_{\{1\}}(x)dx=\overline{\xi(1)}=1.$$
  Korzystając z twierdzenia Fouriera mamy
  $$1=g(1)=\int\overline{\langle 1, \xi\rangle}\Hat{g}(\xi)d\xi=\int 1d\xi=d\xi(\Hat{G})$$
\end{proof}











\subsection{Dualność Pontrjagina}

Cel: $G\cong \Hat{\Hat{G}}$ przez $x\mapsto \langle x, -\rangle$.

\begin{theorem}{Bochnera}{}
  Jeśli $\phi\in\mathcal{P}(G)$ to istnieje jedyna dodatnia miara $\mu\in M(\Hat{G})$ taka, że $\phi=\phi_\mu$.

  \color{red}Gal robi dla niedodatnich tylko tak ogólniej?
\end{theorem}

\begin{proof}
  BSO $\phi(1)=1$. Definiujemy formę Hermitowską
  $$\langle f,g\rangle_\phi=\int\phi(g^*\ast f)$$
  z nierówności Schwarza:
  $$\left|\int\phi(g^*\ast f)\right|^2\leq \int \phi(f^*\ast f)\int \phi(g^*\ast g)$$
  dla $f,g\in L^1G$.
\end{proof}

\begin{zadanie}{3}{10}
  Niech $f\in L^1A\cap L^2A$. Niech $\Hat{\eta}$ będzie współczynnikiem diagonalnym $f$ (z twierdzenia Bochnera). Zrozum napis (wiadomo co to całka z miary; całka z funkcji to całka względem miary Haara):
  $$\|f\|_2^2=\taH{\eta}(1)=\int\eta=\int\Hat{\taH{\eta}}=\int|\Hat{f}|^2=\|\Hat{f}\|^2_2.$$
\end{zadanie}

\begin{proof}
  Co robi Folland? Folland mówi
  $$\|f\|_2^2=\int f^2=(f\ast f^*)(1)=\int\Hat{f\ast f^*}(\xi)d\xi=\int |\Hat{f}(\xi)|^2d\xi,$$
  gdzie 
  $$f\ast f^*(x)=\int f(y)f^*(y^{-1}x)dy=\int f(y)\Delta(x^{-1}y)\overline{f(x^{-1}y)}=\int f(y)\overline{f(x^{-1}y)}$$
  bo mamy grupę abelową.

  $\color{green}\|f\|_2^2=\taH{\eta}(1)$

  Moje podejrzenia: bierzemy $f$ jak wyżej i robimy z tego na siłę funkcję dodatnio określoną, czyli $f\ast f^*$. Z twierdzenia Bochnera mamy $\eta\in M(\Hat{G})$ takie, że 
  $$(f\ast f^*)(x)=\int\langle x, \xi\rangle d\eta(\xi)(=\Hat{\eta}(x)).$$

  Czyli jak u Follanda dla miary $\mu\in M(\Hat{G})$ robimy funkcję z $C_b(G)$ przez $\mu\mapsto \phi_\mu$ i odwracamy to przez $\phi\mapsto \mu_\phi$, to u Gala to odwracanie to jest $\taH{\phi}=\mu_\phi$. 
  
  $\color{green}\taH{\eta}(1)=\int \eta$

  Z twierdzenia Fouriera o inwersji/odwracaniu (?) mam
  $$\taH{\eta}(x)=\int \langle x, \xi\rangle d\eta(\xi)$$
  no a jeśli $x=1$ to mam po prostu po prawej $\int d\eta(\xi)$.

  $\color{green}\int\eta=\int\Hat{\taH{\eta}}$

  Z twierdzenia Fouriera o odwracaniu (?) wiemy, że $d\eta(\xi)=\Hat{\taH{\eta}}(\xi)d\xi$. 

  $\color{green}\int\Hat{\taH{\eta}}=\int |\Hat{f}|^2$
  
  Po pierwsze z definicji $f^*$ i transformaty Foueriera mam
  \begin{align*}
    \Hat{f^*}(\xi)&=\int \overline{\langle x, \xi\rangle}f^*(x)dx=\int\overline{\langle x,\xi\rangle f(x^{-1})}dx=\\ 
                  &=[x^{-1}=u,\;dx=du]=\int \overline{\langle u^{-1}, \xi\rangle f(u)}du=\\ 
                  &=\int\langle u,\xi\rangle \overline{f(u)}du=\overline{\int \overline{\langle x, \xi\rangle}f(x)dx}=\\ 
                  &=\overline{\Hat{f}(\xi)}
  \end{align*}

  Tutaj wiem, że $\taH{\eta}=f\ast f^*$, czyli
  $$\int \Hat{\taH{\eta}}(\xi)d\xi=\int \Hat{f\ast f^*}(\xi)d\xi=\int \Hat{f}(\xi)\Hat{f^*}(\xi)d\xi$$
  i z wyliczenia wyżej dostaję $\int \Hat{f}(\xi)\overline{\Hat{f}(\xi)}d\xi=\int |\Hat{f}(\xi)|^2d\xi$.
\end{proof}

\begin{zadanie}{3}{11}
  Wywnioskuj twierdzenie Reigleya-Parsevala-Plancherela: transformata Fouriera jednoznacznie rozszerza się do unitarnej izometrii $L^2A\to L^2\Hat{A}$.
\end{zadanie}

\begin{proof}
  Zbiór $L^1 G\cap L^2G$ jest gęsty w $L^2G$, bo Stone-Weierstrass lub pewnie można delikatniej. 

  Zadanie wcześniej pokazuje, że $f\mapsto \Hat{f}$ jest izometrią w $L^2$ normie na $L^1G\cap L^2 G$. Czyli rozszerza się do izometrii $L^2G\to L^2\Hat{G}$ (która zachowuje granice ciągów -> chyba w ten sposób bym ja konstruowała wogóle).

  Czyli mamy w ręku izometrię $F:L^2G\to L^2\Hat{G}$, czyli będzie $F^*F=Id$. Teraz że $FF^*=Id$, czyli że jest unitarna -> Folland pokazuje surjektywność. Tutaj mamy elegancki iloczyn skalarny, czyli żeby pokazać surjektywność $F$ wystarczy pokazać, że $F^\perp=0$. W takim razie bierzemy $\phi$ ortogonalne do wszystkich $\Hat{f}$ dla $f\in L^1G\cap L^2G$. To oznacza, że 
  $$0=\langle \phi, \Hat{f}\rangle_{L^2}=\int \phi\overline{\Hat{f}}$$
  w szczególności gdy pod $f$ podstawimy $L_\omega f$ dla dowolnego $\omega$ to dostajemy
  $$0=\int \phi\overline{\Hat{L_\omega f}}=\int \phi(\xi)\overline{\overline{\langle\omega, \xi\rangle}\Hat{f}(\xi)}=\int \phi(\xi)\langle \omega, \xi\rangle \overline{\Hat{f}(\xi)}.$$
  Dalej, skoro $\phi, \Hat{f}\in L^2\Hat{G}$, to $\phi\Hat{f}\in L^1\Hat{G}$, czyli $\phi(\xi)\overline{\Hat{f}(\xi)}d\xi$ definiuje miarę na $\Hat{G}$. Ale $\mu\to \phi_\mu$ dla $\phi_\mu(x)=\int \langle x, \xi\rangle d\mu(\xi)$ jest injekcją, więc skoro tutaj dla dowolnego $f$ mamy zero, to znaczy że miara $\phi\overline{\Hat{f}}d\xi$ jest felerna, czyli $\phi\overline{\Hat{f}}=0$ dla wszystkich $f$, czyli $\phi=0$.
\end{proof}



\end{document}
